\section*{Tóm tắt}
\addcontentsline{toc}{section}{Tóm tắt}

Báo cáo phân tích chuỗi thời gian lưu lượng dòng chảy sông Colorado trong giai đoạn
2010--2023, sử dụng \DatasetN{} ngày quan sát từ \DateMin{} đến \DateMax{}.
Dữ liệu được kiểm tra tính đầy đủ, tính duy nhất và tính liên tục trước khi sử dụng.
Báo cáo trả lời ba nhóm câu hỏi: (i) mô tả xu hướng, mùa vụ và cấu trúc tự tương quan
của chuỗi; (ii) xác định các ngày và các đợt lưu lượng vượt ngưỡng cực trị Q95;
(iii) xây dựng và đánh giá mô hình dự báo trực tiếp 7 ngày, so sánh với các mô hình
đơn giản, đồng thời công bố kết quả trên một website tĩnh.

Kết quả chính cho thấy chuỗi có tính dai dẳng rất mạnh: hệ số tự tương quan bậc 7
là \AcfSeven, bậc 30 là \AcfThirty, trong khi tương quan với độ trễ một năm gần như
bằng không (\AcfYear). Phân rã STL cho thấy thành phần xu hướng chiếm ưu
thế so với thành phần mùa vụ; biên độ trung bình theo tháng chỉ khoảng
\MonthlyAmplitude{} m\textsuperscript{3}/s, trong khi biên độ theo năm lên tới
\YearlyAmplitude{} m\textsuperscript{3}/s. Ngưỡng Q95 của toàn bộ chuỗi là
\QHigh{} m\textsuperscript{3}/s, tương ứng \NDaysHigh{} ngày vượt ngưỡng
(\ShareHigh{}\%) và \NEvents{} đợt liên tiếp. Đợt dài nhất kéo dài
\LongestEvent{} ngày; lưu lượng lớn nhất quan sát được là \QMax{}
m\textsuperscript{3}/s.

Với bài toán dự báo 7 ngày, mô hình LightGBM trực tiếp đạt RMSE \LgbmMeanRmse{}
và NSE \LgbmMeanNse{} trên tập Test, gần tương đương mô hình Naive
(RMSE \NaiveMeanRmse{}, NSE \NaiveMeanNse{}) và tốt hơn rất nhiều so với Seasonal
Naive 365 (RMSE \SeasonalMeanRmse{}). Khoảng cách nhỏ giữa LightGBM và Naive cho
thấy trong chuỗi này, giá trị gần nhất đã chứa phần lớn thông tin dự báo ngắn hạn;
mô hình học máy không tạo ra cải thiện vượt trội. Báo cáo phân tích rõ nguyên nhân,
các yếu tố ảnh hưởng theo từng horizon, mức độ quan trọng của đặc trưng và những
giới hạn của kết quả.

Toàn bộ mã nguồn, dữ liệu đã xử lý, bảng, hình và báo cáo được công bố tại
\url{https://github.com/huyduong5/colorado-river-flow-atlas}. Website kết quả có
địa chỉ \url{https://huyduong5.github.io/colorado-river-flow-atlas/}. Báo cáo không
sử dụng GRU/LSTM do môi trường không cài PyTorch; mọi số liệu trong báo cáo đều
được sinh lại từ pipeline, không có kết quả mô phỏng hoặc số liệu tự điền.

\vspace{0.5em}
\noindent\textbf{Từ khóa:} chuỗi thời gian, lưu lượng dòng chảy, Q95, STL,
LightGBM, dự báo trực tiếp, RMSE, NSE, tái lập.

\newpage
\section*{Nội dung bài tập và phạm vi thực hiện}
\addcontentsline{toc}{section}{Nội dung bài tập và phạm vi thực hiện}

Bài tập yêu cầu phân tích một chuỗi thời gian thực tế và trình bày kết quả trong
một báo cáo có thể tái lập. Ba nhóm câu hỏi được triển khai như sau:

\begin{enumerate}[leftmargin=1.6em]
  \item \textbf{Phân tích chuỗi:} mô tả phân bố, xu hướng dài hạn, tính mùa vụ và
  cấu trúc tự tương quan của lưu lượng dòng chảy; kiểm tra các giả định về tính
  dừng và độ dài mùa.
  \item \textbf{Phát hiện cực trị:} xác định ngưỡng Q95, đếm số ngày vượt ngưỡng,
  nhận diện các đợt liên tiếp, mô tả phân bố theo năm và theo tháng; kiểm tra độ
  nhạy của kết quả khi dùng ngưỡng Z-score tĩnh.
  \item \textbf{Dự báo 7 ngày:} xây dựng mô hình dự báo trực tiếp cho từng horizon,
  so sánh Naive, Seasonal Naive 365 và LightGBM; đánh giá bằng RMSE, MAE và NSE
  trên tập Test; phân tích lỗi theo horizon, theo năm và theo mức lưu lượng; công
  bố bảng, hình và website tương tác.
\end{enumerate}

Phạm vi dữ liệu là 14 năm, từ \DateMin{} đến \DateMax{}. Tập Test gồm hai
năm cuối (2022--2023) và được giữ kín trong toàn bộ quá trình chọn mô hình. Đây là
đánh giá backtest trên dữ liệu lịch sử, không phải hệ thống cảnh báo lũ vận hành.
Báo cáo không đưa thêm biến ngoại sinh như mưa, nhiệt độ, tuyết tan hoặc vận hành
hồ chứa vì bộ dữ liệu đi kèm không có các cột này.

\newpage
\tableofcontents
\newpage
\listoftables
\newpage
\listoffigures

\newpage
\section*{Danh mục từ viết tắt}
\addcontentsline{toc}{section}{Danh mục từ viết tắt}

\begin{table}[H]
\centering
\small
\begin{tabularx}{\linewidth}{@{}lX@{}}
\toprule
\textbf{Từ viết tắt} & \textbf{Diễn giải} \\
\midrule
ACF & Hàm tự tương quan (Autocorrelation Function) \\
MAE & Sai số tuyệt đối trung bình (Mean Absolute Error) \\
NSE & Hệ số hiệu quả Nash--Sutcliffe (Nash--Sutcliffe Efficiency) \\
Q95 & Phân vị 95\% của lưu lượng trên toàn bộ chuỗi \\
RMSE & Căn bậc hai sai số bình phương trung bình (Root Mean Squared Error) \\
STL & Phân rã xu hướng -- mùa vụ bằng LOESS (Seasonal-Trend decomposition using LOESS) \\
LightGBM & Mô hình cây tăng cường gradient của Microsoft \\
CI/CD & Tích hợp liên tục và triển khai liên tục \\
SSG & Trang tĩnh được tạo trước (Static Site Generation) \\
\bottomrule
\end{tabularx}
\caption{Danh mục từ viết tắt dùng trong báo cáo.}
\end{table}
